\section*{Hoofdstuk \ref{ch:g8:ice-on-earth} --- IJs op Aarde}

\begin{solution}{exo:g8:ice-on-earth:1}
(a) Het accumulatiegebied. (b) Het ablatiegebied. (c) De tong. (d) Een
(eind)\omterm{def:g8:ice-on-earth:glacial-landforms}{morene}.
\end{solution}

\begin{solution}{exo:g8:ice-on-earth:2}
(a) Niet waar: een \omterm{def:g8:ice-on-earth:glacier}{gletsjer} bestaat uit \omterm{def:g4:clouds-rain-snow:precipitation}{sneeuw} die op het land is gevallen
en in ijs is veranderd; bevroren zeewater is \omterm{def:g8:ice-on-earth:sea-ice}{zee-ijs}. (b) Waar. (c) Niet
waar: zijn ijs stroomt nog steeds naar beneden; de tong wijkt terug omdat
het ijs sneller smelt dan het aanstroomt. (d) Waar.
\end{solution}

\begin{solution}{exo:g8:ice-on-earth:3}
(a) Een \omterm{def:g8:ice-on-earth:glacial-landforms}{U-dal}. (b) Een zwerfsteen. (c) Een \omterm{def:g8:ice-on-earth:ice-sheet}{ijsplaat}. (d) \omterm{def:g8:ice-on-earth:permafrost}{Permafrost}.
\end{solution}

\begin{solution}{exo:g8:ice-on-earth:4}
(a) Antarctica en Groenland. (b) Antarctica, met ongeveer \qty{58}{m}
(Groenland: ongeveer \qty{7.4}{m}). (c) Ongeveer een tiende (10~\%).
\end{solution}

\begin{solution}{exo:g8:ice-on-earth:5}
(a) Ongeveer \num{7.05}~miljoen~\unit{km^2} in 1979 en \num{3.57}~miljoen~\unit{km^2}
in 2012. (b) $7.05 - 3.57 = 3.48$: ongeveer \num{3.5}~miljoen~\unit{km^2}.
(c) $3.48 \div 7.05 \approx 0.49$: ongeveer de helft. (d) Een enkel jaar
hangt af van het \omterm{def:g1:watching-the-weather:weather}{weer} van die zomer: 2012 ligt ver onder de lijn. De
trendlijn gebruikt alle jaren, dus verandert één ongewone zomer haar maar
heel weinig.
\end{solution}

\begin{solution}{exo:g8:ice-on-earth:6}
(a) $2004 - 1941 = 63$ jaar. (b) $12000 \div 63 \approx 190.5$: ongeveer
\qty{190}{m} per jaar. (c) $2021 - 1984 = 37$ jaar; $3200 \div 37 \approx
86$: ongeveer \qty{86}{m} per jaar. (d) De Muirgletsjer, meer dan twee
keer zo snel.
\end{solution}

\begin{solution}{exo:g8:ice-on-earth:7}
In het \emph{accumulatiegebied}, hoog op de \omterm{def:g1:landscapes-around-us:landform}{berg}, valt elk jaar meer
\omterm{def:g4:clouds-rain-snow:precipitation}{sneeuw} dan er smelt. Nieuwe \omterm{def:g4:clouds-rain-snow:precipitation}{sneeuw} bedekt de oude \omterm{def:g4:clouds-rain-snow:precipitation}{sneeuw} en drukt die
samen; na ongeveer een jaar wordt die \emph{firn}, en dieper en veel later
massief ijs. Onder zijn eigen gewicht \emph{stroomt} het ijs langzaam het
\omterm{def:g1:landscapes-around-us:landform}{dal} af. Lager, in het \emph{ablatiegebied}, smelt het sneller dan er
\omterm{def:g4:clouds-rain-snow:precipitation}{sneeuw} valt; het ijs dat de tong bereikt, is ijs dat hoog op de \omterm{def:g1:landscapes-around-us:landform}{berg} is
ontstaan, vaak lang geleden.
\end{solution}

\begin{solution}{exo:g8:ice-on-earth:8}
(a) In glas B: het gesmolten ijs loopt in het water en doet het stijgen; in
glas A blijft het peil waar het was. (b) Glas A is de \omterm{def:g8:ice-on-earth:ice-sheet}{ijsberg} (ijs dat al
in de \omterm{def:g2:water-on-earth:ocean}{zee} drijft); glas B is de \omterm{def:g8:ice-on-earth:ice-sheet}{ijskap} (ijs op het land). (c) Een \omterm{def:g8:ice-on-earth:ice-sheet}{ijsplaat}
houdt de \omterm{def:g8:ice-on-earth:glacier}{gletsjers} erachter tegen. Als ze uiteenvalt, gaan deze \omterm{def:g8:ice-on-earth:glacier}{gletsjers}
sneller stromen en voeren ze meer landijs naar de \omterm{def:g2:water-on-earth:ocean}{zee}, en dat doet de \omterm{def:g2:water-on-earth:ocean}{zee}
stijgen.
\end{solution}

\begin{solution}{exo:g8:ice-on-earth:9}
(a) Een \omterm{def:g8:ice-on-earth:glacial-landforms}{U-dal}, uitgeslepen door een \omterm{def:g8:ice-on-earth:glacier}{gletsjer}. (b) Een \omterm{def:g8:ice-on-earth:glacier}{gletsjer} heeft het
van ver boven in het \omterm{def:g1:landscapes-around-us:landform}{dal} meegevoerd, waar graniet voorkomt, en het
achtergelaten toen het ijs smolt: het is een zwerfsteen. (c) Een
\omterm{def:g8:ice-on-earth:glacial-landforms}{zijmorene}: haar hoogte laat zien hoe dik de \omterm{def:g8:ice-on-earth:glacier}{gletsjer} ooit was; het ijs
reikte minstens zo hoog op de dalwand.
\end{solution}

\begin{solution}{exo:g8:ice-on-earth:10}
(a) $7.4 \div 10 = \qty{0.74}{m}$. (b) $58 \div 100 = \qty{0.58}{m}$.
(c) $58 \div 7.4 \approx 7.8$: ongeveer 8 keer meer.
\end{solution}

\begin{solution}{exo:g8:ice-on-earth:11}
(a) $7.21 \times 10^6 = \qty{7210000}{km^2}$ en $4.65 \times 10^6 =
\qty{4650000}{km^2}$. (b) $7.21 - 4.65 = \num{2.56}\text{ miljoen }\unit{km^2}$, dat
is \qty{2560000}{km^2}. (c) $2.56 \div 3 \approx 0.85$: ongeveer
\num{0.85}~miljoen~\unit{km^2} per decennium, iets meer dan de
\num{0.76}~miljoen~\unit{km^2} van de trendlijn; elk decenniumgemiddelde
draagt nog de schommelingen van zijn eigen tien zomers in zich, maar beide
manieren om de grafiek te lezen tonen hetzelfde grote verlies. (d) De
uitgestrektheid kan niet onder nul zakken, en een trendlijn beschrijft
alleen de jaren die zijn gemeten: wat er daarna gebeurt, hangt af van hoe
het \omterm{def:g5:weather-and-climate:climate}{klimaat} van het Noordpoolgebied verandert, dus de lijn is geen
voorspelling.
\end{solution}

\begin{solution}{exo:g8:ice-on-earth:12}
(a) $120 \div 1 = 120$: de \omterm{def:g1:landscapes-around-us:water}{kust} lag ongeveer \qty{120}{km} verder in \omterm{def:g2:water-on-earth:ocean}{zee}.
(b) Het was uit de \omterm{def:g2:water-on-earth:ocean}{oceaan} verdampt, als \omterm{def:g4:clouds-rain-snow:precipitation}{sneeuw} op de \omterm{def:g3:continents-and-oceans:continent}{continenten} gevallen
en daar gebleven als het ijs van de grote \omterm{def:g8:ice-on-earth:ice-sheet}{ijskappen}. (c) Een \omterm{def:g8:ice-on-earth:glacial-landforms}{eindmorene}
ontstaat aan het uiteinde van een \omterm{def:g8:ice-on-earth:glacier}{gletsjer}. Een \omterm{def:g8:ice-on-earth:glacial-landforms}{eindmorene} op een \omterm{def:g1:landscapes-around-us:landform}{vlakte}
ver van alle \omterm{def:g1:landscapes-around-us:landform}{bergen} laat zien dat een \omterm{def:g8:ice-on-earth:ice-sheet}{ijskap} die \omterm{def:g1:landscapes-around-us:landform}{vlakte} ooit heeft
bereikt, in een tijd waarin ijs veel meer land bedekte dan nu.
\end{solution}

\begin{solution}{pb:g8:ice-on-earth:1}
\textbf{1.} $1080 - 0 = \qty{1080}{m}$.

\textbf{2.} 1960--1970: \qty{180}{m}; 1970--1980: \qty{150}{m};
1980--1990: \qty{190}{m}; 1990--2000: \qty{180}{m}; 2000--2010:
\qty{180}{m}; 2010--2020: \qty{200}{m}. Het minst in 1970--1980, het meest
in 2010--2020.

\textbf{3.} Nee: een \omterm{def:g8:ice-on-earth:glacier}{gletsjer} die zich terugtrekt, stroomt nog steeds naar
beneden. Zijn tong wijkt terug omdat het ijs sneller smelt dan het
aankomt.

\textbf{4.} Een \omterm{def:g8:ice-on-earth:glacial-landforms}{eindmorene}: het ijs voerde puin het \omterm{def:g1:landscapes-around-us:landform}{dal} af en liet het
vallen bij zijn tong, waar het ijs smolt, terwijl de tong daar stond.

\textbf{5.} $1080 \div 6 = \qty{180}{m}$ per decennium.

\textbf{6.} $180 \div 10 = \qty{18}{m}$ per jaar.

\textbf{7.} $190 \div 18 \approx 10.6$: ongeveer 11 keer sneller.

\textbf{8.} Minder: hij verloor in zijn ablatiegebied meer ijs dan hij in
zijn accumulatiegebied won, dus week zijn tong terug.

\textbf{9.} Een U-vorm: een brede bodem met steile wanden, omdat de
\omterm{def:g8:ice-on-earth:glacier}{gletsjer} het \omterm{def:g1:landscapes-around-us:landform}{dal} verbreedde en verdiepte terwijl hij het vulde.

\textbf{10.} $1980 - 1080 = \qty{900}{m}$.

\textbf{11.} $900 \div 18 = 50$ jaar: in \textbf{2070}.

\textbf{12.} De snelheid veranderde van decennium tot decennium (van 150
tot \qty{200}{m}), en de toekomstige terugtocht hangt af van de sneeuwval
en de zomertemperaturen van de komende jaren: een snelheid uit het
verleden is geen voorspelling.
\end{solution}
